\honest{Beginnings: An Introduction}{Rob Bensinger}{2015}

This last book is less a conclusion than a call to action, so I will point you to further reading.

The book's definition of normative rationality, in terms of probability theory and decision theory, is standard in cognitive science. For introductions, read Baron's \textsc{Thinking and Deciding} and the \textsc{Oxford Handbook of Thinking and Reasoning}.

The book's philosophy is ``more controversial.'' Yudkowsky argues that a rational agent should take one box in Newcomb's problem, ``a minority position among working decision theorists.'' \nb{The PhilPapers surveys bear this out. In 2020, one box drew 22.7 per cent of decision theorists' answers other than ``Other.''} Gary Drescher's \textsc{Good and Real} reaches many of the same conclusions on philosophy of science and decision theory independently, and so serves as a book-length treatment of this book's core philosophy. Talbott surveys the views within Bayesian epistemology, among them Jaynes's position that not all priors are equally reasonable. Like E. T. Jaynes, Yudkowsky wants a rule for choosing priors, not only for revising beliefs, which aligns Yudkowsky with AI researchers such as Marcus Hutter, who hope to understand general AI through a better theory of ideal reasoning. For efforts to naturalize theories of knowledge, see Feldman.

Some frequentists say Bayesian probabilities are subjective. Kruschke and Yudkowsky reply that frequentism is more subjective, because its probabilities depend on the experimenter's intentions. \nb{This is the stopping-rule argument. The frequentist answer to it is discussed in the notes on ``Beautiful Probability.''} The philosophical dispute is not the same as the choice of statistical methods; both kinds are useful when used correctly. Bayesian tools have become cheaper since the 1980s, and their informativeness and generality are more widely appreciated, which has produced ``Bayesian revolutions'' in many sciences. Frequentist methods remain more popular, and in some contexts are still better. Kruschke's \textsc{Doing Bayesian Data Analysis} is an accessible introduction.

Training in statistics, and in some other fields such as psychology, improves reasoning outside the classroom; classes in formal logic and informal fallacies ``have not proven similarly useful.'' \nb{The two studies cited are, by their titles and the summaries the editor could read, about statistical training. The claims about psychology and about logic classes could not be checked.}

The book has three sequences. ``Yudkowsky's Coming of Age'' gives a last example of how irrational belief works, taken from the author's own intellectual history. ``Challenging the Difficult'' asks what it takes to solve a truly hard problem, including demands beyond getting beliefs right. ``The Craft and the Community'' is about rationality groups and group rationality. It raises these questions. Can rationality be learned and taught, and if so, how much improvement is possible? How can we be sure a rationality intervention has a real effect, and find its cause? What community norms would help? Can we work together on large problems without giving up freedom of thought and conduct? Above all, what is missing? The next primers should replace this text, improve its style, ``test its prescriptions,'' add to its content and go in new directions.

Yudkowsky wrote these essays because of his own philosophical mistakes and his difficulties in AI theory, but they proved useful to a much wider audience. The posts inspired the growth of \textsc{Less Wrong}, a community interested in cognitive science, computer science and philosophy. Its writers helped seed effective altruism, an effort to find the charities and causes that do the most good. The posts also sparked the Center for Applied Rationality, a nonprofit that tries to turn the science of rationality into usable techniques for self-improvement. I do not know what comes next. The art of applied rationality is ``a new and half-formed thing''; there are few rationalists and much left undone. ``May you serve your purpose well.''
